\section{Proofs for \texorpdfstring{Section~\ref{sec:coh}}{coherence}}\label{app:coh}
% src: L3 proofs (Sections 1-3, 4.2, 4.4, 5), L7-P8, L6-P7, L4-P19, B5 Sections 3 and 5 proofs, critic V7 and V8 re-derivations; ring-rate lower bound, ring wall counting, transfer-matrix formula, repeated circulation and relay recursions added here

This appendix proves the results of \cref{sec:coh} in the order in which they appear. It also states the secondary results to which the section points, \cref{lem:coh-one-message,lem:coh-rank-one,lem:coh-c1-value} and \cref{conj:coh-walls}. Notation is that of the section.

\subsection{Components and fidelity}\label{app:coh-fidelity}

\begin{lemma}[One message]\label{lem:coh-one-message}
Let $j$'s data be jointly Gaussian with $\eta\sim\Ncal(0,P_0)$, let $m=\E[\eta\mid\text{data of }j]$ have posterior covariance $P$, so that $\Cov(m)=P_0-P$, and let $y=H_jm+\nu$ with $\nu\sim\Ncal(0,R_c)$ independent and $R_c\succ0$. Then $y=H_j(P_0-P)P_0^{-1}\eta+\nu'$ with $\nu'$ independent of $\eta$, $\Cov(\nu')=H_j(P-PP_0^{-1}P)H_j^\top+R_c$, and
\begin{equation}\label{eq:coh-one-message}
I(\eta;y)=\tfrac12\ln\det\bigl(H_j(P_0-P)H_j^\top+R_c\bigr)-\tfrac12\ln\det\bigl(H_j(P-PP_0^{-1}P)H_j^\top+R_c\bigr).
\end{equation}
This information is at most $\min\{I(\eta;m),I(m;y)\}$, where $I(m;y)=\frac12\sum_k\ln(1+x_k)$ over the at most $r_c$ nonzero eigenvalues $x_k$ of $R_c^{-1/2}H_j(P_0-P)H_j^\top R_c^{-1/2}$.
\end{lemma}
% src: L3-P1

\begin{proof}
The Bayes error $\eta-m$ is orthogonal to every function of $j$'s data, so $\Cov(\eta,m)=\Cov(m)=P_0-P$. The regression of $m$ on $\eta$ is therefore $(P_0-P)P_0^{-1}\eta$, and its residual $\tilde m$ is independent of $\eta$, the pair being jointly Gaussian and uncorrelated, with covariance $(P_0-P)-(P_0-P)P_0^{-1}(P_0-P)=P-PP_0^{-1}P$. Substituting $m=(P_0-P)P_0^{-1}\eta+\tilde m$ into $y=H_jm+\nu$ gives the representation with $\nu'=H_j\tilde m+\nu$. For jointly Gaussian vectors $I(\eta;y)=\frac12\ln\det\Cov(y)-\frac12\ln\det\Cov(y\mid\eta)$, with $\Cov(y)=H_j(P_0-P)H_j^\top+R_c$ and $\Cov(y\mid\eta)=\Cov(\nu')$, which is \eqref{eq:coh-one-message}. The variables $\eta$, $m$, $y$ form a Markov chain, so data processing gives both upper bounds, and $I(m;y)=\frac12\ln\det(I+R_c^{-1/2}H_j\Cov(m)H_j^\top R_c^{-1/2})$ is the stated sum, with at most $\rank H_j\le r_c$ nonzero terms.
\end{proof}

\begin{proof}[Proof of \cref{prop:coh-sdpi}]
We use one cited fact. For the Ornstein--Uhlenbeck kernel $\mathsf O_\tau(x,\cdot)=\Ncal(e^{-\tau}x,(1-e^{-2\tau})I)$ in any dimension and every probability measure $Q$ with $\KL(Q\Vert\Ncal(0,I))<\infty$,
\begin{equation}\label{eq:coh-ou-decay}
\KL(Q\mathsf O_\tau\Vert\Ncal(0,I))\le e^{-2\tau}\KL(Q\Vert\Ncal(0,I)).
\end{equation}
This is the decay of relative entropy along the Ornstein--Uhlenbeck semigroup, equivalent to the Gaussian logarithmic Sobolev inequality with constant one \citep{gross1975logarithmic}; we use it in the form of \citet[Theorem~5.2.1]{bakry2014analysis} as restated by \citet[Section~2.3]{ledoux2015stein}.

Since $M_{XY}M_{XY}^\top=I-\Cov(Y)^{-1/2}R_{\mathsf K}\Cov(Y)^{-1/2}\prec I$, every canonical correlation is below one. Take a singular value decomposition $M_{XY}=V_0D_0W_0^\top$ and put $X'=W_0^\top\Cov(X)^{-1/2}X$ and $Y'=V_0^\top\Cov(Y)^{-1/2}Y$. Both are standard Gaussian vectors, $\Cov(Y',X')=D_0$ and $Y'\mid X'\sim\Ncal(D_0X',I-D_0D_0^\top)$: with $\sigma_k=\sigma_k(M_{XY})$ and independent standard $Z_k$, $Y'_k=\sigma_kX'_k+(1-\sigma_k^2)^{1/2}Z_k$ for $k\le\min\{\dim X,\dim Y\}$ and $Y'_k=Z_k$ beyond. Relative entropy is invariant under these bijections, so it suffices to treat $(X',Y')$. If $\sigma_1=0$ the output is independent of the input and every claim is trivial. Otherwise put $\tau_k=-\ln\sigma_k$, with $\tau_k=\infty$ when $\sigma_k=0$ or $k>\min\{\dim X,\dim Y\}$. The channel $X'\mapsto Y'$ is the composition of three kernels: the isotropic kernel $\mathsf O_{\tau_1}$; the product kernel that applies to coordinate $k$ the one-dimensional Ornstein--Uhlenbeck kernel at time $\tau_k-\tau_1\ge0$, time $\infty$ meaning replacement by independent standard noise; and the map that keeps the first $\min\{\dim X,\dim Y\}$ coordinates and, if $\dim Y>\dim X$, appends independent standard coordinates. The composite is the channel because Ornstein--Uhlenbeck kernels compose additively in time: $e^{-\tau'}(e^{-\tau}x+(1-e^{-2\tau})^{1/2}Z_1)+(1-e^{-2\tau'})^{1/2}Z_2$ has the law $\mathsf O_{\tau+\tau'}(x,\cdot)$. The second and third kernels do not increase relative entropy, the second by data processing and the third because dropping coordinates is data processing and appending an independent common factor leaves relative entropy unchanged. The first contracts it by $e^{-2\tau_1}=\sigma_1^2$, by \eqref{eq:coh-ou-decay}. Hence $\KL(Q_X\mathsf K\Vert P_X\mathsf K)\le\sigma_1^2\KL(Q_X\Vert P_X)$.

For attainment let $w_1$ be the top right singular vector of $M_{XY}$ and $\alpha\ne0$ a scalar, and shift the mean along it:
\[ Q_X=\Ncal\bigl(\alpha\Cov(X)^{1/2}w_1,\Cov(X)\bigr),\qquad Q_X\mathsf K=\Ncal\bigl(\alpha B\Cov(X)^{1/2}w_1,\Cov(Y)\bigr). \]
Then $\KL(Q_X\Vert P_X)=\alpha^2/2$ and $\KL(Q_X\mathsf K\Vert P_X\mathsf K)=\frac{\alpha^2}2\Vert M_{XY}w_1\Vert^2=\frac{\alpha^2}2\sigma_1^2$.

For the mutual-information form, the Markov property gives $P_{Y\mid U}=P_{X\mid U}\mathsf K$ and $P_Y=P_X\mathsf K$, so $I(U;Y)=\E_U\KL(P_{Y\mid U}\Vert P_Y)\le\sigma_1^2\E_U\KL(P_{X\mid U}\Vert P_X)=\sigma_1^2I(U;X)$. For $U=X'_1+Z'$ with $Z'\sim\Ncal(0,1/\alpha)$, $\alpha>0$, independent of everything else, $I(U;X)=\frac12\ln(1+\alpha)$ and $I(U;Y)=-\frac12\ln(1-\sigma_1^2\alpha/(1+\alpha))$, and their ratio tends to $\sigma_1^2$ as $\alpha\to0$, so the constant is best possible.
\end{proof}

\begin{proof}[Proof of \cref{prop:coh-relay}]
Write $\tilde X_k=\Cov(X_k)^{-1/2}X_k$. For a centered Gaussian Markov chain
\[ \E[\tilde X_k\mid\tilde X_{k-1},\dots,\tilde X_0]=M_k\tilde X_{k-1}, \]
so the tower property gives $\E[\tilde X_L\mid\tilde X_0]=M_L\cdots M_1\tilde X_0$ and $\Cov(\tilde X_L,\tilde X_0)=M_L\cdots M_1$, which is (a), the operator norm being submultiplicative. The pair $(X_0,X_L)$ is jointly Gaussian, so $X_L$ is a linear function of $X_0$ plus independent Gaussian noise, whose covariance $\Cov(X_L\mid X_0)\succeq\Cov(X_L\mid X_{L-1})\succ0$ by the Markov property. \Cref{prop:coh-sdpi} applies with $M_{X_0X_L}=M_L\cdots M_1$, and $U$, $X_0$, $X_L$ form a Markov chain, which gives (b). For scalars (a) is the product of the correlations, and $I(X_0;X_L)=-\frac12\ln(1-\Corr(X_0,X_L)^2)$ for jointly Gaussian scalars, which is (c).
\end{proof}

Relay conventions. Let a source sense a static scalar world $\eta\sim\Ncal(0,P_0)$ through noise of variance $R^s$, let every hop add independent channel noise of variance $R_c$, and let relays not sense. Under posterior-mean forwarding the message after $u$ hops is $y_u=m_{u-1}+n_u$, with $m_0=\E[\eta\mid s]$ and $m_u=\E[\eta\mid y_u]$. Writing $v_u=\Var(m_u)=\Cov(\eta,m_u)$, the message has $\Cov(\eta,y_u)=v_{u-1}$ and $\Var(y_u)=v_{u-1}+R_c$, so $v_u=v_{u-1}^2/(v_{u-1}+R_c)\le v_{u-1}^2/R_c$ and $I(\eta;y_u)=-\frac12\ln(1-v_u/P_0)$; once $v_u<R_c$ the information decays doubly exponentially. Under unit-variance forwarding every relay sends its content scaled to unit variance, so each hop multiplies the squared correlation with $\eta$ by $1/(1+R_c)$, starting from $P_0/(P_0+R^s)$, and the information decays exponentially. Under unbiased forwarding every relay sends what it received, so after $u$ hops the message is $\eta$ plus noise of variance $R^s+uR_c$, and $I=\frac12\ln(1+P_0/(R^s+uR_c))$ decays like $P_0/(2uR_c)$. In a persistent world $\eta_{t+1}=F\eta_t+w_t$ with $|F|<1$, let relay $k$ update $m_k(t+1)=F[(1-g)m_k(t)+gy_{k-1}(t)]$, where $y_{k-1}(t)$ is $m_{k-1}(t)$ plus independent channel noise. The innovation $w_t$ and the channel noise are independent of $m(t)$, so in stationarity $c_k=\Cov(\eta_t,m_k(t))$ satisfies $c_k=F^2[(1-g)c_k+gc_{k-1}]$, that is, $c_k=\psi c_{k-1}$ with $\psi=F^2g/(1-F^2(1-g))$. Since $g\le1$, $\psi\le F^2$, and $\psi\to1$ as $F\to1$ at fixed $g>0$.

\begin{proof}[Proof of \cref{prop:coh-info-decay}]
(a) Under (Lin) the schedule is independent of $U$ and of the noises, and $I(U;Z)\le I(U;Z,(W(t))_t)=I(U;Z\mid(W(t))_t)$, so it suffices to prove the bound for a fixed schedule. Write $\Psi(t,t')=W(t-1)\cdots W(t')$, with $\Psi(t,t)=I$ (the transition product written $\Phi(t,t')$ in \cref{sec:sca}, the letter $\Phi$ being reserved here for the code potential). Unrolling the recursion, $e_i(t)=\mathfrak aU+Z_e$ with $\mathfrak a=\sum_{t'\le s}\Psi_{ij}(t,t'+1)g_j(t')$, where $Z_e$ collects the initial errors and the other noise components; it is Gaussian by (G) and independent of $U$. The row-sum norm is submultiplicative, so $|\Psi_{ij}(t,t'+1)|\le\lambda^{t-1-t'}$, and under (Loc) the entry vanishes when $t-1-t'<d(i,j)$, which is the light cone of \cref{thm:sca-lightcone}. The exponents $t-1-t'$ that contribute are therefore at least $n_0=\max\{d(i,j),t-1-s\}$, and $|\mathfrak a|\le\bar g\sum_{n\ge n_0}\lambda^n=\bar g\lambda^{n_0}/(1-\lambda)$. The mean is $m_i(t)=\eta_t+e^{\mathrm{com}}_i(t)+e_i(t)=\mathfrak aU+Z_m$ with $Z_m=\eta_t+e^{\mathrm{com}}_i(t)+Z_e$, where the common mode $e^{\mathrm{com}}_i(t)$ of \cref{lem:fw-common-mode} is a linear function of the world path; $Z_m$ is Gaussian and independent of $U$, because $U$ is independent of the world, the initial errors and the other noise components jointly. Both $Z_e$ and $Z_m$ contain, with coefficient one, the component of $\zeta_i(t-1)$ that is independent of everything else, so their variances are at least $q_{\min}$. For $Z$ Gaussian and independent of $U$, $I(U;\mathfrak aU+Z)=h(\mathfrak aU+Z)-h(Z)\le\frac12\ln(1+\mathfrak a^2\Var(U)/\Var(Z))\le\mathfrak a^2\Var(U)/(2\Var(Z))$, where $h(\cdot)$ is differential entropy and a Gaussian maximizes it at fixed variance. Combining the bounds gives (a).

(b) On the ring $W$ is symmetric and nonnegative with $W\one=\lambda\one$, and $\lambda<1$ gives the stationary process. Place $i$ at $0$ and $j$ at $d=d(i,j)$. The local noises injected at $j$ enter $m_i(t)$ only through $X_j=\sum_{n\ge0}(W^n)_{0d}\zeta_j(t-1-n)$, and $m_i(t)=X_j+Y_j$ with $Y_j$ Gaussian and independent of the noise sequence at $j$, so $m_i(t)$ depends on that sequence only through $X_j$. Hence
\[ I=I(X_j;X_j+Y_j)=-\frac12\ln\Bigl(1-\frac{qG(d)}{\Var(m_i(t))}\Bigr),\qquad G(d)=\sum_{n\ge0}\bigl((W^n)_{0d}\bigr)^2, \]
and since $\Var(m_i(t))$ does not depend on $d$, the rate of $I$ is that of $G(d)$. The entries $(W^n)_{0d}$ are nonnegative, and their sum over $n$ is the resolvent
\[ \sum_{n\ge0}(W^n)_{0d}=\frac1{2\pi}\int_{-\pi}^{\pi}\frac{e^{\mathrm iud}}{1-c-2a\cos u}\mathrm{d}u=\frac{e^{-md}}{\sqrt{(1-c)^2-4a^2}},\qquad\cosh m=\frac{1-c}{2a}=1+\frac{\varepsilon}{2a}, \]
by the standard evaluation of the integral by residues; $(1-c)^2>4a^2$ because $\lambda<1$. Upper bound: a sum of squares of nonnegative terms is at most the square of their sum, so $G(d)\le e^{-2md}/((1-c)^2-4a^2)$. Lower bound: fix $A_0>m/\ln(1/\lambda)$. Since $(W^n)_{0d}\le\lambda^n$, the terms with $n>A_0d$ sum to at most $\lambda^{A_0d}/(1-\lambda)$, which is $o(e^{-md})$, so for large $d$ the terms with $n\le A_0d$ sum to at least half the resolvent, and by the Cauchy--Schwarz inequality
\[ G(d)\ge\frac{1}{A_0d+1}\Bigl(\sum_{n\le A_0d}(W^n)_{0d}\Bigr)^2\ge\frac{e^{-2md}}{4(A_0d+1)((1-c)^2-4a^2)}. \]
The two bounds give $-\ln G(d)/d\to2m$, with a prefactor between a multiple of $1/d$ and a constant. Since $I=qG(d)(1+o(1))/(2\Var(m_i(t)))$ as $d\to\infty$, the same holds for $I$.
\end{proof}

\subsection{Calibration and withdrawal}\label{app:coh-calibration}

\begin{proof}[Proof of \cref{prop:coh-c0-blind}]
(a) The expected free energy of a policy is computed from the receiver's own predictive density $\int p_i(o\mid\eta,\pi)q_i(\eta)\mathrm{d}\eta$ of outcomes together with its prior preferences. Under $\mathsf{C0}$ the likelihood of translated messages does not contain $\Upsilon_{ij}$, so the expected free energy depends on the mistranslations only through the belief $q_i$, which they influence through past data. (b) In the static two-group model the believed information of an attention set $\mathcal J$ is $\frac K2\ln(\Lambda_i(\mathcal J)/\Lambda_i(\varnothing))$ with $\Lambda_i(\mathcal J)=\sigma_\eta^{-2}+1/R^s_i+\sum_{j\in\mathcal J}g_j^2/\hat r^0_{ij}$, where $\hat r^0_{ij}$ is the fixed assigned noise, for example $R_c+g_j^2R^s_j$. This involves neither the data nor $\Upsilon_{ij}$, so the policy posterior, proportional to $E_i(\pi)\exp(-\gamma G_i(\pi))$, does not depend on the mistranslations at any $\gamma$, and two senders with equal assigned noise contribute equal terms.
\end{proof}

\begin{proof}[Proof of \cref{thm:coh-withdrawal}]
(a) Write $\Upsilon=\Upsilon_{ij}$. The translated message is $\tilde y_{ij}=\Upsilon g_j(\eta+v_j)+\nu_{ij}$, so $\tilde y_{ij}-g_j\eta=(\Upsilon-I)g_j\eta+g_j\Upsilon v_j+\nu_{ij}$, a sum of independent terms. With $\Upsilon$ orthogonal, $\E\Vert\Upsilon v_j\Vert^2=KR^s_j$, and under $\mathsf{C1g}$ the population value $\hat r_{ij}=K^{-1}\E\Vert\tilde y_{ij}-g_j\eta\Vert^2$ equals $R_c+g_j^2R^s_j+g_j^2\sigma_\eta^2\Vert\Upsilon-I\Vert_F^2/K$. Under $\mathsf{C1i}$ the receiver's estimate is $\hat\eta_i=g_is_i$, whose error has variance $g_iR^s_i$ per coordinate and is uncorrelated with $\hat\eta_i$, so its model assigns the innovation $\tilde y_{ij}-g_j\hat\eta_i$ the variance $g_j^2g_iR^s_i+\hat r_{ij}$ per coordinate. The actual innovation is $g_j(\Upsilon-g_iI)\eta+g_j\Upsilon v_j+\nu_{ij}-g_jg_iv_i$, with second moment per coordinate $g_j^2\sigma_\eta^2\Vert\Upsilon-g_iI\Vert_F^2/K+g_j^2R^s_j+R_c+g_j^2g_i^2R^s_i$. For orthogonal $\Upsilon$, $\Vert\Upsilon-g_iI\Vert_F^2=K(1-g_i)^2+g_i\Vert\Upsilon-I\Vert_F^2$, both sides being $K(1+g_i^2)-2g_i\tr\Upsilon$, and $(1-g_i)^2\sigma_\eta^2+g_i^2R^s_i=g_iR^s_i$ by the definition of $g_i$. Matching the model value to the actual one gives $\hat r_{ij}=R_c+g_j^2R^s_j+g_ig_j^2\sigma_\eta^2\Vert\Upsilon-I\Vert_F^2/K$.

(b) Under $\mathsf{C1}$ the receiver's model treats attended messages as independent isotropic observations of $g_j\eta$ with noise variance $\hat r_{ij}$, so the believed information of an attention set $\mathcal J$ is $\frac K2\ln(\Lambda_i(\mathcal J)/\Lambda_i(\varnothing))$ with $\Lambda_i(\mathcal J)=\sigma_\eta^{-2}+1/R^s_i+\sum_{j\in\mathcal J}\varpi_{ij}$. This is strictly increasing in $\sum_{j\in\mathcal J}\varpi_{ij}$, and every $\varpi_{ij}$ is positive, so the maximizers over $|\mathcal J|\le\kappa$ consist of $\kappa$ largest calibrated precisions, ties broken arbitrarily, or of all senders if fewer are available.

(c) If $\varpi_{ib}<\varpi^{\mathrm{in}}_{i(\kappa)}$, at least $\kappa$ in-group senders have calibrated precision above $\varpi_{ib}$, so by (b) the cross-group sender $b$ is not attended; with $\varpi_{ib}=g_b^2/(R_c+g_b^2R^s_b+\hat r^{\mathrm{mt}}_{ib})$ this condition is \eqref{eq:coh-withdrawal}. If $\varpi_{ib}>\varpi^{\mathrm{in}}_{i(\kappa)}$ for some cross-group sender $b$, every attention set of $\kappa$ in-group senders contains a sender of precision at most $\varpi^{\mathrm{in}}_{i(\kappa)}$, and replacing it by $b$ raises the sum of precisions, so by (b) the receiver attends at least one cross-group sender. When no cross-group precision equals $\varpi^{\mathrm{in}}_{i(\kappa)}$ the two cases are exhaustive, which gives the equivalence. With homogeneous sensing noise all in-group precisions equal $g^2/(R_c+g^2R^s)$, the right side of \eqref{eq:coh-withdrawal} vanishes, and $\hat r^{\mathrm{mt}}_{ib}>0$ whenever $\Upsilon_{ib}\ne I$, because $g_i>0$.
\end{proof}

\subsection{Consistency}\label{app:coh-consistency}

For a walk $p$ write $\Theta(p)$ for the ordered product of the translations along it, with inverses on steps against an edge.

\begin{proof}[Proof of \cref{thm:coh-holonomy}]
(a) We restate the proof of \citet{dennis2026magent} for linear translations. If $\mu$ is in consensus, composing the edge relations around a directed closed walk $\mathfrak c$ at $i_0$ gives $\mu_{i_0}=\Hol(\mathfrak c)\mu_{i_0}$. Conversely, let $v$ be fixed by all such holonomies and set $\mu_i=\Theta(p)v$ along a directed path $p$ from $i_0$ to $i$. A second path $p'$ gives the same value: strong connectivity supplies a directed path $p''$ from $i$ back to $i_0$, the products $\Theta(p'')\Theta(p)$ and $\Theta(p'')\Theta(p')$ are holonomies of directed closed walks at $i_0$ and both fix $v$, and $\Theta(p'')$ is invertible. Extending a directed path to $j$ by the edge $j\to i$ gives a directed path to $i$, so $\mu_i=\Theta_{ij}\mu_j$. The map $\mu\mapsto\mu_{i_0}$ is linear, and it is injective because every $\mu_i$ is determined by $\mu_{i_0}$ along a path.

(b) Let $\Theta_{\mathcal T}(i_0\to i)=\Theta(p)$ for the path $p$ from $i_0$ to $i$ in $\mathcal T$. Each edge relation can be solved for either endpoint, so by induction along tree paths every consensus configuration satisfies $\mu_i=\Theta_{\mathcal T}(i_0\to i)\mu_{i_0}$, and $\mu\mapsto\mu_{i_0}$ is injective. For a non-tree edge $e=(j\to i)$ the relation $\mu_i=\Theta_{ij}\mu_j$ then reads $\Hol_ev=v$ with $v=\mu_{i_0}$ and $\Hol_e=\Theta_{\mathcal T}(i_0\to i)^{-1}\Theta_{ij}\Theta_{\mathcal T}(i_0\to j)$, the holonomy of the closed walk described in the statement. Conversely, for $v$ fixed by every $\Hol_e$ the configuration $\mu_i=\Theta_{\mathcal T}(i_0\to i)v$ satisfies the tree relations by construction and the non-tree relations by the fixed-point condition.

(c) The first statement restates (b); on strongly connected graphs both fixed spaces are images of $\mathrm{Cons}(\Theta)$ under the same injective map, so they coincide. If every $\Hol_e$ is the identity, put $g_i=\Theta_{\mathcal T}(i_0\to i)$; then $\Theta_{ij}=g_ig_j^{-1}$ on tree edges by construction and on non-tree edges because $g_i^{-1}\Theta_{ij}g_j=\Hol_e=I$. Conversely, the holonomy of a coboundary telescopes to the identity around every closed walk. Since $\dim\mathrm{Cons}(\Theta)=\dim\bigcap_e\mathrm{Fix}(\Hol_e)$, full consensus holds exactly when every $\Hol_e$ is the identity.
\end{proof}

\begin{proof}[Proof of \cref{prop:coh-flat-true}]
Along a closed walk $i_0,i_1,\dots,i_L=i_0$ every step contributes $U_{i_u}U_{i_{u-1}}^{-1}$, whether it runs along an edge, contributing $\Theta_{i_ui_{u-1}}$, or against one, contributing $\Theta_{i_{u-1}i_u}^{-1}=(U_{i_{u-1}}U_{i_u}^{-1})^{-1}$. The product telescopes to $U_{i_0}U_{i_0}^{-1}=I$, and for general encodings the cancellations $\mathcal E_{i_u}^{-1}\circ\mathcal E_{i_u}$ give the identity map. By \cref{thm:coh-holonomy}(c) the consensus space has dimension $K$; the configurations $\mu_i=U_ix$ satisfy every edge relation and form a $K$-dimensional space, so they are all of it, and by \cref{prop:coh-frustration}(a) the least eigenvalue of the frustration form is zero. The last statement is the contrapositive: if the translations are the true ones and every agent encodes one true world invertibly, the least frustration vanishes.
\end{proof}

Learned holonomy. In the tangent learning model of \citet{dennis2026magent}, receiver $i$ learns a generator $\omega_e\in\mathfrak{gl}(K)$ of its translation $\Theta_{ij}=e^{\omega_e}$ on each edge $e=(j\to i)$ from grounded messages, whose law depends on $\omega_e$ only through $T_e\omega_e$, with $T_e$ the stacked public codes of the modalities read on $e$. The model is linear in the generators: a message on $e$ has mean $T_e(I-\omega_e)\bar y$ for a referent $\bar y$, the first-order form of the grounded message model at $\Theta_{ij}=e^{\omega_e}$, with Gaussian channel noise, and the prior on the generators is Gaussian. The convergence statement used here is part (ii) of the theorem on learning with a loop prior in the chapter on communication codes and learned alignment of \citet{dennis2026magent}. Its hypotheses are a positive-definite prior precision, referents drawn independently with a finite positive-definite second moment, observed source labels, and $\lceil\upsilon_en\rceil$ independent exchanges on every edge $e$ with rates $\upsilon_e>0$; its conclusion concerns the mean of the exact Gaussian posterior, so no step-size schedule enters, and holds almost surely as $n\to\infty$ by the strong law of large numbers. With an isotropic anchor, no loop prior and grounded messages on every edge, that posterior mean converges to $(I-\Pi_N)\omega^\star$, where $\omega^\star$ are the true generators and $\Pi_N$ is the orthogonal projection onto the generators that no message distinguishes from zero, those with $T_e\omega_e=0$ on every edge; on edge $e$ it acts by left multiplication with the orthogonal projector $\Pi_e$ onto $\ker T_e$. For private encodings near a common frame, taken as the identity, the true generators are, to first order in $a_i=\log U_i$, the coboundary $\omega^\star_e=a_i-a_j$, which sums to zero around every directed closed walk. By the Baker--Campbell--Hausdorff formula the learned holonomy of a directed closed walk is therefore $I-\sum_e\Pi_e(a_i-a_j)$ to first order, the sum running over the edges $e=(j\to i)$ of the walk. Its first-order part vanishes when $\Pi_e$ is the same on every edge of the walk, as when all its edges are read through one modality. To first order in the frame error it is not detectable from messages: the learned generators differ from $\omega^\star$ by $-\Pi_e\omega^\star_e$ on each edge, and $T_e\Pi_e=0$, so in the tangent model every message has the same law under the learned generators as under the true ones, whose translations are flat (\cref{prop:coh-flat-true}). Beyond the tangent model the message means under $e^{(I-\Pi_e)\omega^\star_e}$ and $e^{\omega^\star_e}$ differ at second order, by $\frac12T_e\omega^\star_e\Pi_e\omega^\star_e\bar y$ to leading order, so the statement does not extend to the exponentiated translations at that order. The first-order formula is checked in \cref{chk:coh-holonomy}.

\begin{proof}[Proof of \cref{prop:coh-frustration}]
(a) By the definition of $\mathcal B$, $\Vert\mathcal B\mu\Vert^2=\mathfrak F(\mu)$. Since every $\varpi_{ij}>0$ and $M_i\succ0$, $\mathcal B\mu=0$ exactly when $\mu_i=\Theta_{ij}\mu_j$ on every edge, so $\ker\mathcal L=\ker\mathcal B=\mathrm{Cons}(\Theta)$. By the Rayleigh principle $\lambda_{\min}(\mathcal L)=\min_{\Vert\mu\Vert=1}\mathfrak F(\mu)$, which is positive exactly when $\mathrm{Cons}(\Theta)=0$, that is, by \cref{thm:coh-holonomy}(b), when the $\Hol_e$ have no common nonzero fixed vector.

(b) An orthogonal change of internal coordinates $\mu_k\mapsto g_k\mu_k$ replaces $\Theta_{k+1,k}$ by $g_{k+1}\Theta_{k+1,k}g_k^{-1}$ and conjugates $\mathcal L$ by an orthogonal block-diagonal matrix, so it preserves the spectrum. Choosing $g_1=I$ and $g_{k+1}=g_k\Theta_{k+1,k}^{-1}$ for $k=1,\dots,N-1$ makes every translation the identity except the one on the edge from $N$ to $1$, which becomes the holonomy $\Hol$ based at agent $1$. Over $\mathbb C$ the orthogonal matrix $\Hol$ is diagonalized by a unitary matrix, applied at every agent, with eigenvalues $e^{\mathrm i\theta_r}$, $\theta_r\in(-\pi,\pi]$, and the frustration decouples into $K$ scalar rings in which one edge carries the phase $e^{\mathrm i\theta_r}$ and the others carry one. The diagonal unitary change of variables $\mu'_k=e^{\mathrm i(k-1)\theta_r/N}\mu_k$ gives every edge the phase $e^{\mathrm i\theta_r/N}$, and the resulting operator is translation invariant on $\Z/N$. Its eigenvectors are the waves $\mu'_k=e^{2\pi\mathrm iuk/N}$, $u=0,\dots,N-1$, whose residual on every edge has squared modulus $|e^{2\pi\mathrm iu/N}-e^{\mathrm i\theta_r/N}|^2=2-2\cos((2\pi u-\theta_r)/N)$ times $|\mu'_k|^2$, which gives the spectrum. For $u=0$ the angle $(2\pi u-\theta_r)/N$ has modulus $|\theta_r|/N$, and for $1\le u\le N-1$ it lies in $[\pi/N,2\pi-\pi/N]$, at distance at least $\pi/N\ge|\theta_r|/N$ from $2\pi\Z$, so the least eigenvalue is $2-2\cos(\phi/N)=4\sin^2(\phi/(2N))$.

(c) Lower bound. Take a simple cycle $\mathfrak c\in\mathcal C_i$ traversed as $i=i_0,i_1,\dots,i_L=i$, let $\Theta^{(u)}$ be the transport of step $u$ as in \cref{def:coh-components}(c) and $r_u=\mu_{i_u}-\Theta^{(u)}\mu_{i_{u-1}}$ its residual. Unrolling gives $(I-\Hol^{(i)}(\mathfrak c))\mu_i=\sum_{u=1}^L\Theta^{(L)}\cdots\Theta^{(u+1)}r_u$. The transports are orthogonal, so the residual of a step against an edge has the norm of that edge's residual and every product of transports has norm one; the Cauchy--Schwarz inequality gives $\Vert(I-\Hol^{(i)}(\mathfrak c))\mu_i\Vert^2\le L\sum_u\Vert r_u\Vert^2\le L\mathfrak F(\mu)/\varpi_{\min}$, since a simple cycle uses each edge once. Dividing by $L=|\mathfrak c|$ and averaging over $\mathcal C_i$ gives $\mu_i^\top\mathcal D_i\mu_i\le\mathfrak F(\mu)/\varpi_{\min}$ for every $i$. Choosing the agent with the largest $\Vert\mu_i\Vert$, for which $\Vert\mu_i\Vert^2\ge\Vert\mu\Vert^2/N$, gives $\mathfrak F(\mu)\ge\varpi_{\min}N^{-1}\min_i\lambda_{\min}(\mathcal D_i)\Vert\mu\Vert^2$. Upper bound. For a unit vector $v$ the configuration $\mu_i=\Theta_{\mathcal T}(i_0\to i)v$ has $\Vert\mu\Vert^2=N$, zero residual on tree edges, and residual $\Theta_{\mathcal T}(i_0\to i)(I-\Hol_e)v$, of norm $\Vert(I-\Hol_e)v\Vert$, on a non-tree edge $e=(j\to i)$; minimizing $\mathfrak F(\mu)/\Vert\mu\Vert^2$ over $v$ gives the bound. Single cycle. Take $\mathcal C_i$ to be the cycle itself for every $i$. The holonomies based at different agents are conjugate, and $(I-\Hol)^\top(I-\Hol)=2I-\Hol-\Hol^\top$ has the eigenvalues $2-2\cos\theta_r$, because $\Hol$ is normal, so the lower bound is $4\sin^2(\phi/2)/N^2$. With $\sin(\phi/(2N))\le\phi/(2N)$ and $\sin(\phi/2)\ge\phi/\pi$ on $[0,\pi]$, the exact value $4\sin^2(\phi/(2N))$ of (b) is at most $\phi^2/N^2\le(\pi^2/4)\cdot4\sin^2(\phi/2)/N^2$.
\end{proof}

\begin{proof}[Proof of \cref{prop:coh-coded}]
(a) $\delta_{\mathfrak S}$ maps $\R^{NK}$ to $\R^{|\mathcal E_{\mathrm a}|r_c}$, so its kernel has dimension at least $NK-|\mathcal E_{\mathrm a}|r_c$. (b) The holonomy of the cycle is the rotation by $0.3+0.5+\pi/2-0.8=\pi/2$, which fixes no nonzero vector, so $\mathrm{Cons}(\Theta)=0$ by \cref{thm:coh-holonomy}. Number the agents so that the edges are $1\to2$, $2\to3$ and $3\to1$, and let $\alpha_1$, $\alpha_2$, $\alpha_3$ be the cosines of the rotation angles of $\Theta_{21}$, $\Theta_{32}$, $\Theta_{13}$, which lie in $(0,1)$ because the angles lie in $(0,\pi/2)$. For a rotation by $\phi$, $(1\ 0)$ times its inverse is $(\cos\phi\ \ \sin\phi)$, so the three rows of $\delta_{\mathfrak S}$, restricted to the first coordinates of $\mu_1$, $\mu_2$ and $\mu_3$, form the matrix with rows $(-1,\alpha_1,0)$, $(0,-1,\alpha_2)$ and $(\alpha_3,0,-1)$, whose determinant is $\alpha_1\alpha_2\alpha_3-1\ne0$. Hence $\delta_{\mathfrak S}$ has rank $3$ and $\dim H^0(\mathfrak S)=6-3=3$.
\end{proof}

\begin{proof}[Proof of \cref{prop:coh-frustration-grounding}]
$\mathcal L=\mathcal B^\top\mathcal B$ is symmetric and positive semidefinite, and with orthogonal translations $\mu^\top\mathcal L\mu=\sum_e\Vert\mu_i-\Theta_{ij}\mu_j\Vert^2\le2\sum_e(\Vert\mu_i\Vert^2+\Vert\mu_j\Vert^2)\le2\Delta_{\mathrm a}\Vert\mu\Vert^2$, every agent lying on at most $\Delta_{\mathrm a}$ edges. So the spectrum of $\mathcal L$ lies in $[0,2\Delta_{\mathrm a}]$ and that of $I-\alpha\mathcal L$ in $[0,1]$, and the norm of $I-\alpha\mathcal L$ is its largest eigenvalue, $1-\alpha\lambda_{\min}(\mathcal L)$. Multiplying by $1-\varepsilon$ gives the contraction factor and $\varepsilon_{\mathrm{cal}}=1-(1-\varepsilon)(1-\alpha\lambda_{\min}(\mathcal L))$. The value on the cycle is \cref{prop:coh-frustration}(b).
\end{proof}

\begin{proof}[Proof of \cref{prop:coh-attended-frustration}]
With orthogonal frames each translation is $\Theta_{ij}=U_i\Upsilon_{ij}U_j^{-1}$, so along a closed walk the intermediate frames cancel and $\Hol(\mathfrak c)=U_{i_0}\Upsilon^{(L)}\cdots\Upsilon^{(1)}U_{i_0}^{-1}$, where $\Upsilon^{(u)}$ is the mistranslation of step $u$, inverted on steps against an edge; inversion preserves $\Vert\Upsilon-I\Vert$ for orthogonal $\Upsilon$. Telescoping, $\Upsilon^{(L)}\cdots\Upsilon^{(1)}-I=\sum_u(\Upsilon^{(L)}\cdots\Upsilon^{(u+1)})(\Upsilon^{(u)}-I)$, and orthogonal factors and orthogonal conjugation preserve norms, so $\Vert\Hol(\mathfrak c)-I\Vert\le|\mathfrak c|\delta_{\mathrm{att}}$. An attended cycle graph of $N$ agents can be oriented as a directed cycle without changing its frustration, since an edge term $\Vert\mu_i-\Theta_{ij}\mu_j\Vert$ equals $\Vert\mu_j-\Theta_{ij}^{-1}\mu_i\Vert$ for orthogonal $\Theta_{ij}$. Let $\phi\in[0,\pi]$ be the smallest modulus of an eigen-angle of its holonomy. By \cref{prop:coh-frustration}(b), $\lambda_{\min}(\mathcal L)=4\sin^2(\phi/(2N))\le(\phi/N)^2$. Also $2\sin(\phi/2)\le\Vert\Hol-I\Vert\le N\delta_{\mathrm{att}}$, the norm being $2\sin$ of half the largest modulus of an eigen-angle, and $\phi\le\pi\sin(\phi/2)$ on $[0,\pi]$, so $\phi/N\le\pi\delta_{\mathrm{att}}/2$.
\end{proof}

\subsection{Amplification}\label{app:coh-amplification}

\begin{proof}[Proof of \cref{thm:coh-amplification}]
(a) For the edge $e=(i\to j)$ put
\[ \tilde A_e=P_j^{-1/2}A_{ji}P_i^{1/2}=\bigl[P_j^{1/2}H_i^\top S_{ji}^{-1/2}\bigr]\bigl[S_{ji}^{-1/2}H_iP_i^{1/2}\bigr]. \]
Around the cycle the factors $P_k^{\pm1/2}$ cancel, so $P_1^{-1/2}A_{\mathfrak c}P_1^{1/2}=\tilde A_{(N\to1)}\cdots\tilde A_{(1\to2)}$, and $\rho(A_{\mathfrak c})\le\prod_e\Vert\tilde A_e\Vert$, because the spectral radius is invariant under similarity and bounded by any submultiplicative norm. The first bracket has squared norm $\lambda_{\max}(S_{ji}^{-1/2}X_{ji}S_{ji}^{-1/2})=\mathrm{snr}_e/(1+\mathrm{snr}_e)$. When $\hat R_i\succ0$, writing $X_{ji}=\hat R_i^{1/2}Y\hat R_i^{1/2}$ shows that $S_{ji}^{-1}X_{ji}$ is similar to $(I+Y)^{-1}Y$, whose eigenvalues are $x/(1+x)$ for the eigenvalues $x$ of $Y$, which identifies $\mathrm{snr}_e$. The second bracket has squared norm $\lambda_{\max}(S_{ji}^{-1/2}X_{ii}S_{ji}^{-1/2})=\mathfrak r_e$. Since $\mathrm{snr}_e/(1+\mathrm{snr}_e)\le1$, \eqref{eq:coh-edge-ratio} follows.

(b) For noiseless codes $H_i\Kg_{ji}=H_iP_jH_i^\top S_{ji}^{-1}=I$, so every edge map is idempotent. A receiver $j$ whose mean is perturbed by $x_j$ and that hears a sender perturbed by $x_i$ moves to $(I-A_{ji})x_j+A_{ji}x_i$. Let agent 1 be perturbed by $y$ at the start of a circulation and let each agent update when it hears its predecessor. In the first circulation agents $2,\dots,N$ hold no perturbation; in later ones agent $k\ge2$ holds $A_{k,k-1}$ applied to its predecessor's previous perturbation, and its own term drops out because $(I-A_{k,k-1})A_{k,k-1}=0$. In either case agent $k$ ends the circulation with $A_{k,k-1}\cdots A_{21}y$, and agent 1 ends it with $(I-A_{1N})y+A_{\mathfrak c}y=\bar A_{\mathfrak c}y$, so after $n$ circulations agent 1 holds $\bar A_{\mathfrak c}^nx$ when it started with $x$. Since $A_{1N}$ is idempotent and $A_{\mathfrak c}=A_{1N}A_{\mathfrak c}$, in the decomposition $\mathrm{ran}(A_{1N})\oplus\ker(A_{1N})$ the map $\bar A_{\mathfrak c}$ is block upper triangular, with diagonal blocks $A_{\mathfrak c}$ restricted to $\mathrm{ran}(A_{1N})$ and the identity on $\ker(A_{1N})$, which is nonzero because the code read by agent 1 has rank below $K$. The eigenvectors of $A_{\mathfrak c}$ with nonzero eigenvalues lie in $\mathrm{ran}(A_{1N})$, so $\rho(\bar A_{\mathfrak c})=\max\{\rho(A_{\mathfrak c}),1\}$, and $\Vert\bar A_{\mathfrak c}^n\Vert$ grows geometrically exactly when this exceeds one.
\end{proof}

\begin{proof}[Proof of \cref{prop:coh-no-amplification}]
(a) With $P_i=P_j$, $X_{ii}=X_{ji}$, so $\mathfrak r_e=\mathrm{snr}_e/(1+\mathrm{snr}_e)$, and \eqref{eq:coh-edge-ratio} gives $\rho(A_{\mathfrak c})\le\prod_e\mathrm{snr}_e/(1+\mathrm{snr}_e)$, which is at most one, and below one when every $\hat R_i\succ0$.

(b) $P_i\preceq e^{\operatorname{sp}(P_i\Vert P_j)}P_j$, hence $X_{ii}\preceq e^{\operatorname{sp}(P_i\Vert P_j)}X_{ji}$ and $\mathfrak r_e\le e^{\operatorname{sp}(P_i\Vert P_j)}\mathrm{snr}_e/(1+\mathrm{snr}_e)$. By \eqref{eq:coh-edge-ratio}, $\rho(A_{\mathfrak c})\le\prod_ee^{\operatorname{sp}(P_i\Vert P_j)/2}\mathrm{snr}_e/(1+\mathrm{snr}_e)$, which is at most one under the stated condition. If all priors lie within Thompson distance $\delta_{\mathrm T}$ of $P$, then $P_i\preceq e^{\delta_{\mathrm T}}P\preceq e^{2\delta_{\mathrm T}}P_j$, so $\operatorname{sp}(P_i\Vert P_j)\le2\delta_{\mathrm T}$, and the condition holds when $N\delta_{\mathrm T}\le\sum_e\ln(1+\mathrm{snr}_e^{-1})$.

(c) If $\hat R_i\succeq X_{ii}$, then $S_{ji}\succeq X_{ji}+X_{ii}$. Since $S_{ji}\succ0$, the generalized Rayleigh quotient $\bm w^\top X_{ii}\bm w/\bm w^\top S_{ji}\bm w$ attains its maximum $\mathfrak r_e$ at some $\bm w\ne0$. If $H_i^\top\bm w=0$, then $\mathfrak r_e=0$. Otherwise $\bm w^\top X_{ji}\bm w=\Vert P_j^{1/2}H_i^\top\bm w\Vert^2>0$, so $\bm w^\top X_{ii}\bm w<\bm w^\top(X_{ii}+X_{ji})\bm w\le\bm w^\top S_{ji}\bm w$ and $\mathfrak r_e<1$. By \eqref{eq:coh-edge-ratio}, $\rho(A_{\mathfrak c})<1$. A level-1 receiver assigns $\hat R_i=H_iP_iH_i^\top+R_c\succeq X_{ii}$ when the sender posts its prior covariance $P_i$.

(d) With a common code $H$ and noise $\hat R$, put $X_j=HP_jH^\top$ and $B_j=X_j(X_j+\hat R)^{-1}$; then $HA_{j,j-1}=B_jH$, so $A_{\mathfrak c}=P_1H^\top(X_1+\hat R)^{-1}B_NB_{N-1}\cdots B_2H$. Since the nonzero spectra of $UV$ and $VU$ agree, with $U=P_1H^\top(X_1+\hat R)^{-1}$ and $V=B_N\cdots B_2H$, the nonzero spectrum of $A_{\mathfrak c}$ is that of $B_NB_{N-1}\cdots B_2B_1$. Each $\hat R^{-1/2}B_j\hat R^{1/2}=Y_j(I+Y_j)^{-1}$ is symmetric with eigenvalues in $[0,1)$, so $\rho(A_{\mathfrak c})\le\prod_j\Vert Y_j(I+Y_j)^{-1}\Vert<1$. For a common code $\bm h^\top$ of rank one with positive sender-specific scalar noises, $\bm h^\top A_{j,j-1}=x_j\bm h^\top$ with $x_j=\bm h^\top P_j\bm h/(\bm h^\top P_j\bm h+\hat R_{j-1})\in[0,1)$, and the only nonzero eigenvalue of $A_{\mathfrak c}$ is $\prod_jx_j<1$.
\end{proof}

For codes of rank one the cycle eigenvalue factorizes.

\begin{lemma}[Codes of rank one]\label{lem:coh-rank-one}
Let every sender $i$ on the cycle use a code $\bm h_i^\top$ of rank one with scalar assigned noise $\hat R_i\ge0$, write $\langle\bm x,\bm x'\rangle_P=\bm x^\top P\bm x'$, $\Vert\bm x\Vert_P=\langle\bm x,\bm x\rangle_P^{1/2}$ and $\cos_P$ for the corresponding cosine, $j^-$ for the sender that $j$ reads and $i^+$ for the receiver that reads $i$. The cycle product has rank at most one, so its only eigenvalue that can be nonzero is its trace,
\begin{equation}\label{eq:coh-rank-one}
\tr A_{\mathfrak c}=\prod_j\frac{\langle\bm h_j,\bm h_{j^-}\rangle_{P_j}}{\Vert\bm h_{j^-}\Vert^2_{P_j}+\hat R_{j^-}}=\prod_j\cos_{P_j}(\bm h_j,\bm h_{j^-})\cdot\mathfrak M_{\mathfrak c}\cdot\prod_i\frac{\Vert\bm h_i\Vert^2_{P_{i^+}}}{\Vert\bm h_i\Vert^2_{P_{i^+}}+\hat R_i},\qquad\mathfrak M_{\mathfrak c}=\prod_i\frac{\Vert\bm h_i\Vert_{P_i}}{\Vert\bm h_i\Vert_{P_{i^+}}}.
\end{equation}
The factors of the first and last products have modulus at most one, and $\mathfrak M_{\mathfrak c}=1$ for a common prior and for a common code. Hence $|\tr A_{\mathfrak c}|>1$ requires $\mathfrak M_{\mathfrak c}>1$.
\end{lemma}
% src: B5-P15

\begin{proof}
Each edge map has rank one, $A_{j,j^-}=P_j\bm h_{j^-}\bm h_{j^-}^\top/(\Vert\bm h_{j^-}\Vert^2_{P_j}+\hat R_{j^-})$. A product of rank-one maps $\bm a_k\bm b_k^\top$ around a cycle has rank at most one, and its trace, which is its only possibly nonzero eigenvalue, is the product of the inner products $\bm b_{k+1}^\top\bm a_k$ of consecutive factors. Here these are $\bm h_j^\top P_j\bm h_{j^-}/(\Vert\bm h_{j^-}\Vert^2_{P_j}+\hat R_{j^-})$, which gives the first form of \eqref{eq:coh-rank-one}. Writing each factor as $\cos_{P_j}(\bm h_j,\bm h_{j^-})\cdot\frac{\Vert\bm h_j\Vert_{P_j}}{\Vert\bm h_{j^-}\Vert_{P_j}}\cdot\frac{\Vert\bm h_{j^-}\Vert^2_{P_j}}{\Vert\bm h_{j^-}\Vert^2_{P_j}+\hat R_{j^-}}$ and regrouping the last two products by the sender $i=j^-$, so that $j=i^+$, gives the second form. Cosines and noise factors have modulus at most one. For a common prior every ratio in $\mathfrak M_{\mathfrak c}$ is one, and for a common code $\mathfrak M_{\mathfrak c}=\prod_i\Vert\bm h\Vert_{P_i}/\prod_i\Vert\bm h\Vert_{P_{i^+}}=1$, both products running over all agents of the cycle.
\end{proof}

Details of \cref{ex:coh-amplification}. In (i) the edge maps are $A_{ab}=P_a\bm e_1\bm e_1^\top/(1+R_c)$ and $A_{ba}=P_b\bm e_2\bm e_2^\top/(1+R_c)$, since $(P_a)_{11}=(P_b)_{22}=1$, and \cref{lem:coh-rank-one} gives $\tr A_{\mathfrak c}=(P_a)_{21}(P_b)_{12}/(1+R_c)^2=\alpha_a\alpha_bx/(1+R_c)^2$. In (ii) the same lemma gives $\tr A_{\mathfrak c}=(P_a)_{21}(P_b)_{12}=(1-u)^2(1+2u)^2$, which exceeds one for $0<u<\frac12$ because $(1-u)(1+2u)=1+u(1-2u)$. The generalized eigenvalues $x$ of the pair $(P_a,P)$, the roots of $\det(P_a-xP)=0$, satisfy $(1+2u-x)^2=4u^2x/(1-(1-u)^2)$; their product is $(1+2u)^2$, so they are $(1+2u)e^{\pm\vartheta}$ with $\cosh\vartheta=1+2u/((2-u)(1+2u))$. Hence $\operatorname{dist}_{\mathrm T}(P_a,P)=\ln((1+2u)e^\vartheta)=\ln(1+2u)+\vartheta$, and the same holds for $P_b$ by exchanging the coordinates. As $u\to0$, $\vartheta=\sqrt{2u}(1+O(u))$ and $\ln(1+2u)=O(u)$, which gives the expansion.

\subsection{Codes}\label{app:coh-codes}

\begin{lemma}[Value of matches]\label{lem:coh-c1-value}
The value \eqref{eq:coh-value} is nondecreasing and concave in $u$, also under a capacity. For an edge on which $n_{\mathrm{mis}}$ of several binary features, each carried by one coordinate, are mismatched, $\mathsf{C1g}$ calibration gives the noise $\hat r_0+4n_{\mathrm{mis}}g^2\sigma_\eta^2/K$; the believed value of the edge is convex and decreasing in $n_{\mathrm{mis}}$, and the value lost to mistranslation, also when capped at the value of the edge because the receiver withdraws, is concave, nondecreasing and zero at $n_{\mathrm{mis}}=0$, hence subadditive.
\end{lemma}
% src: B5-L1

\begin{proof}
The argument of the logarithm in \eqref{eq:coh-value} is affine and increasing in $u$, because $\varpi^{=}>\varpi^{\neq}$, and the logarithm is concave and increasing; under a capacity, $u$ enters through the concave nondecreasing map $u\mapsto\min\{u,\kappa\}$, and a concave nondecreasing function of a concave function is concave. For the edge, the receiver models the message as $g\eta$ plus isotropic noise whose calibrated variance is, by \cref{thm:coh-withdrawal}(a) with $\Vert\Upsilon-I\Vert_F^2=4f$ for $f=n_{\mathrm{mis}}$, $\hat r(f)=\hat r_0+4fg^2\sigma_\eta^2/K$ under $\mathsf{C1g}$, and $\mathsf{C1i}$ multiplies the last term by $g_i$. The function $V(f)=\ln(1+A_1/(A_2+A_3f))$ with positive constants $A_1=g^2/\Lambda^{\mathrm{own}}$, $A_2=\hat r_0$ and $A_3=4g^2\sigma_\eta^2/K$ has
\[ V'(f)=-\frac{A_1A_3}{(A_2+A_3f)(A_2+A_3f+A_1)}<0,\qquad V''(f)=\frac{A_1A_3^2(2(A_2+A_3f)+A_1)}{(A_2+A_3f)^2(A_2+A_3f+A_1)^2}>0, \]
so the edge value is convex and decreasing and the loss $V(0)-V(f)$ is concave and increasing with value zero at $f=0$. The minimum of a concave function and a constant is concave, and a concave function $\ell$ on $[0,\infty)$ with $\ell(0)=0$ satisfies $\ell(f+f')\le\ell(f)+\ell(f')$.
\end{proof}

\begin{proof}[Proof of \cref{prop:coh-adoption}]
With uniform habits the policy posterior over the two codes is $q_i(\sigma)\propto\exp(\gamma\mathcal F_i(u_i(\sigma))+\gamma h_i\sigma)$, and $u_i(-1)=\mathsf d_i-u^+_i$, which gives the logistic form. (a) $\mathcal F_i$ is nondecreasing, so the argument of $\operatorname{sig}$ is nondecreasing in $u^+_i$; the map $u^+_i\mapsto\mathsf d_i-u^+_i$ changes the sign of $\mathcal F_i(u^+_i)-\mathcal F_i(\mathsf d_i-u^+_i)$, and $\operatorname{sig}(-z)=1-\operatorname{sig}(z)$. (b) For $\mathcal F_i(u)=\mathrm{const}+Ju$ the bracket is $J(2u^+_i-\mathsf d_i)=J\mathsf d_i(2x-1)$. Put $\beta_0=\gamma J\mathsf d_i$ and $\psi(x)=\operatorname{sig}(\beta_0(2x-1))-x$. Then $\psi(\frac12)=0$, $\psi'(\frac12)=\beta_0/2-1$, $\psi(1)<0$, and $\psi$ is strictly concave on $[\frac12,1]$ because $\operatorname{sig}$ is strictly concave on $[0,\infty)$. If $\beta_0>2$, $\psi$ is positive just to the right of $\frac12$ and has exactly one root $x_+$ in $(\frac12,1)$, with $\psi'(x_+)<0$ by strict concavity; if $\beta_0\le2$, strict concavity gives $\psi(x)<\psi'(\frac12)(x-\frac12)\le0$ on $(\frac12,1]$, so there is no root. Since $\psi(1-x)=-\psi(x)$, the roots are symmetric about $\frac12$. The map $x\mapsto\operatorname{sig}(\beta_0(2x-1))$ is increasing with derivative $1+\psi'$, which is $\beta_0/2>1$ at $\frac12$ and lies in $(0,1)$ at $x_\pm$, so $\frac12$ is unstable and $x_\pm$ are stable. For $\beta_0>2\ln4$, $\psi(0)>0>\psi(\frac14)$, so $x_-<\frac14$ and $x_-=\operatorname{sig}(-\beta_0(1-2x_-))\le e^{-\beta_0/2}$; then $\beta_0x_-\to0$ and $x_-e^{\beta_0}\to1$. (c) At $x=0$ the bracket is $\mathcal F_i(0)-\mathcal F_i(\mathsf d_i)$, and $\operatorname{sig}$ is positive. The argument of $\operatorname{sig}$ is $\gamma(G_i(-1)-G_i(+1))$, so as $\gamma\to\infty$ the probability tends to one, to $\frac12$ or to zero as $G_i(-1)-G_i(+1)=\mathcal F_i(u^+_i)-\mathcal F_i(\mathsf d_i-u^+_i)+2h_i$ is positive, zero or negative. For $h_i=0$ and strictly increasing $\mathcal F_i$ this difference has the sign of $2u^+_i-\mathsf d_i$, which is the majority rule; under a capacity $\kappa<\mathsf d_i$ the value is constant on $u\ge\kappa$, and the difference vanishes whenever $u^+_i\ge\kappa$ and $\mathsf d_i-u^+_i\ge\kappa$, for example at $\mathsf d_i=3$, $\kappa=1$ and $u^+_i=2$.
\end{proof}

\begin{proof}[Proof of \cref{prop:coh-code-potential}]
(a) For configurations $\sigma$ and $\sigma'$ that differ only at agent $i$, the terms of $\Phi_{\mathrm{code}}$ that involve $\sigma_i$ are $\sum_jJ_{ij}\mathbf 1[\sigma_i\ne\sigma_j]-h_i\sigma_i$, by the symmetry of the couplings, so $G_i(\sigma'_i)-G_i(\sigma_i)=\Phi_{\mathrm{code}}(\sigma')-\Phi_{\mathrm{code}}(\sigma)$ and $\Phi_{\mathrm{code}}$ is an exact potential. The logit chain moves from $\sigma$ to $\sigma'$ with probability $N^{-1}e^{-\gamma G_i(\sigma'_i)}/Z_i$ and back with probability $N^{-1}e^{-\gamma G_i(\sigma_i)}/Z_i$, with the same normalizer $Z_i$, which depends only on the other agents' codes. The ratio is $e^{-\gamma(\Phi_{\mathrm{code}}(\sigma')-\Phi_{\mathrm{code}}(\sigma))}$, which is detailed balance for the law proportional to $e^{-\gamma\Phi_{\mathrm{code}}}$. Every single flip has positive probability, so the chain is irreducible on the finite configuration space and this law is its unique stationary law; as $\gamma\to\infty$ it concentrates on the minimizers of $\Phi_{\mathrm{code}}$.

(b) For an agent with two partners $u_i(-\sigma_i)=2-u_i(\sigma_i)$, and $\mathcal F(u)-\mathcal F(2-u)$ takes the values $-(\mathcal F(2)-\mathcal F(0))$, $0$ and $\mathcal F(2)-\mathcal F(0)$ at $u=0,1,2$, so $G_i(-\sigma_i)-G_i(\sigma_i)=(\mathcal F(2)-\mathcal F(0))(u_i(\sigma_i)-1)+2h_i\sigma_i$. In the pairwise game with coupling $J$ the same flip changes the number of mismatched edges at $i$ by $u_i(\sigma_i)-(2-u_i(\sigma_i))$, so the cost changes by $2J(u_i(\sigma_i)-1)+2h_i\sigma_i$. The two agree in every configuration when $J=\frac12(\mathcal F(2)-\mathcal F(0))$, so the two games differ by terms free of each agent's own code, and $\Phi_{\mathrm{code}}$ is an exact potential for the code game.

(c) If an exact potential exists, the cost changes of the deviating agents along any closed path of unilateral deviations telescope in the potential and sum to zero. Take adjacent agents $i$ and $j$; let $p$ of $i$'s other partners use $+1$ and $\mathsf d_i-1-p$ use $-1$, and $p'$ and $\mathsf d_j-1-p'$ of $j$'s other partners likewise. Starting from $\sigma_i=\sigma_j=+1$, let $i$ flip, then $j$, then $i$ back, then $j$ back. With $\nabla\mathcal F(n)=\mathcal F(n+1)-\mathcal F(n)$ the four cost changes are $\mathcal F_i(p+1)-\mathcal F_i(\mathsf d_i-1-p)+2h_i$, then $\mathcal F_j(p')-\mathcal F_j(\mathsf d_j-p')+2h_j$, then $\mathcal F_i(\mathsf d_i-p)-\mathcal F_i(p)-2h_i$, then $\mathcal F_j(\mathsf d_j-1-p')-\mathcal F_j(p'+1)-2h_j$, and their sum is
\[ \nabla\mathcal F_i(p)+\nabla\mathcal F_i(\mathsf d_i-1-p)-\nabla\mathcal F_j(p')-\nabla\mathcal F_j(\mathsf d_j-1-p'), \]
which is $\nabla^{\pm}\mathcal F_i(p)-\nabla^{\pm}\mathcal F_j(p')$ in the notation of (c). Let $i$ and $j$ have $r$ common partners. The pairs that occur are $p=a+b$ and $p'=a+b'$, where $a\in\{0,\dots,r\}$ counts the common partners that use $+1$ and $b\in\{0,\dots,\mathsf d_i-1-r\}$ and $b'\in\{0,\dots,\mathsf d_j-1-r\}$ count the remaining other partners of $i$ and of $j$ that use $+1$; $p$ and $p'$ range independently only when $r=0$. If $i$ and $j$ have the same other partners, then $b=b'=0$ and $p=p'=a$, and the sum vanishes in every configuration exactly when $\nabla^{\pm}\mathcal F_i(u)=\nabla^{\pm}\mathcal F_j(u)$ for every $u$. Otherwise one endpoint, say $i$, has another partner that is not a partner of $j$, so $\mathsf d_i-1-r\ge1$. Fixing $a$ and $b'$ and varying $b$ shows that $\nabla^{\pm}\mathcal F_i$ is constant on $\{a,\dots,a+\mathsf d_i-1-r\}$; the intervals for consecutive $a$ overlap and cover $\{0,\dots,\mathsf d_i-1\}$, so $\nabla^{\pm}\mathcal F_i$ is constant, and $\nabla^{\pm}\mathcal F_j(a+b')$ equals the same constant for all $a$ and $b'$, which exhaust $\{0,\dots,\mathsf d_j-1\}$. Conversely, the conditions of (c) make the sum vanish for every attended edge and every configuration, and for two agents that are not partners the four changes cancel in pairs because neither cost depends on the other's code; these are all the closed paths of length four, which is the sufficiency stated after the proposition. On a complete graph with a common value $\mathcal F$ any two agents have the same other partners, the condition holds trivially, and $u_i(\sigma_i)=n_{\sigma_i}(\sigma)-1$. When $i$ moves from code $c$ to code $-c$ its cost changes by $\mathcal F(n_c-1)-\mathcal F(n_{-c})+2h_ic$, and $\Phi$ of (c) changes by the same amount, since the sum over code $c$ loses its last term $\mathcal F(n_c-1)$ and the sum over code $-c$ gains the term $\mathcal F(n_{-c})$; so $\Phi$ is an exact potential. For values $\mathrm{const}_i+Ju$ with a slope common to all agents, $G_i(\sigma)$ is $J$ times the number of $i$'s partners that use the other code, minus $h_i\sigma$, up to terms free of $\sigma$: a pairwise game. For a quadratic value $\alpha_0+\alpha_1u+\alpha_2u^2$ on a graph in which every agent has $\mathsf d$ partners, $u_i(\sigma)=\frac12(\mathsf d+\sigma\bar\sigma_i)$ with $\bar\sigma_i$ the sum of the codes of $i$'s partners, and since $\sigma^2=1$, $\mathcal F(u_i(\sigma))=\mathrm{const}+\frac12(\alpha_1+\alpha_2\mathsf d)\sigma\bar\sigma_i+\frac14\alpha_2\bar\sigma_i^2$, whose last term does not involve $\sigma$; this is the pairwise game with $J=\alpha_1+\alpha_2\mathsf d$. For the logarithmic value at $\mathsf d_i=3$ without capacity, the quantity at $u=0$ and at $u=1$ differs by the third difference $\mathcal F(3)-3\mathcal F(2)+3\mathcal F(1)-\mathcal F(0)$, which equals the third derivative of $u\mapsto\frac K2\ln(1+(u\varpi^{=}+(3-u)\varpi^{\neq})/\Lambda^{\mathrm{own}})$ at an intermediate point and is positive. With capacity one, $\mathcal F_i(u)=\mathcal F_i(1)$ for $u\ge1$, so at $\mathsf d_i=3$ the quantity is $\mathcal F_i(1)-\mathcal F_i(0)>0$ at $u=0$ and zero at $u=1$. If an agent $i$ with three partners $j$, $k$ and $l$ has the same other partners as each of them, then the partners of $j$ are $i$, $k$ and $l$, and likewise for $k$ and $l$, so these four agents form a complete graph that no edge leaves, which is the whole graph when the graph is connected. In every other connected graph an agent with three partners has a partner with different other partners, on which edge the necessary condition demands a constant $\nabla^{\pm}\mathcal F_i$, and no exact potential exists. At an end of an open chain of at least three agents the end agent has one partner and its neighbor two, so the two do not have the same other partners; with $x_1=\varpi^{=}/\Lambda^{\mathrm{own}}>x_0=\varpi^{\neq}/\Lambda^{\mathrm{own}}$ the two constants are $K\ln\frac{1+x_1}{1+x_0}$ and $\frac K2\ln\frac{1+2x_1}{1+2x_0}$, and $(1+x_1)^2(1+2x_0)-(1+x_0)^2(1+2x_1)=(x_1-x_0)(x_1+x_0+2x_0x_1)>0$, so they differ; with capacity one the neighbor's constant is $\frac K2\ln\frac{1+x_1}{1+x_0}$, half that of the end agent.
\end{proof}

The proof of \cref{thm:coh-domains-1d} rests on three lemmas. On a ring or an open chain of $N$ agents, a wall at $i$ means $\sigma_i\ne\sigma_{i+1}$, a domain is a maximal run of consecutive agents with one code, and the aligned sum of a run of agents in a domain of sign $\sigma_{\mathsf D}$ is $\sigma_{\mathsf D}$ times the sum of their pulls. Let $\sigma^\ast$ minimize $\Phi_{\mathrm{code}}(\sigma)=J\sum_i\mathbf 1[\sigma_i\ne\sigma_{i+1}]-\sum_ih_i\sigma_i$, the first sum running over the edges of the ring or chain.

\begin{lemma}[Structure of the minimizer]\label{lem:coh-structure}
% src: L3-L3 (extended to rings)
(a) A domain with walls on both sides has aligned sum at least $J$. (b) A proper initial or final run of such a domain has aligned sum at least $0$. (c) Every run inside any domain has aligned sum at least $-J$.
\end{lemma}

\begin{proof}
Flipping a run of agents inside a domain of sign $\sigma_{\mathsf D}$ changes the pull term by $2\sigma_{\mathsf D}$ times the run's sum of pulls and changes the number of walls; optimality of $\sigma^\ast$ makes the total change nonnegative. Flipping a whole domain with walls on both sides removes two walls, which gives (a). Flipping a proper initial or final run of it moves a wall, which gives (b). Flipping a run strictly inside a domain adds two walls, which gives aligned sum at least $-J$; a proper run that touches an end of an open chain adds one wall, which gives at least $-J/2$; flipping a whole end domain of a chain removes one wall, which gives at least $J/2$; and in a configuration without walls every proper run adds one or two walls while the whole population adds none, which gives at least $-J$ and $0$. With (a) and (b) this gives (c) in every case.
\end{proof}

Let $X_0=0$ and $X_i=h_1+\cdots+h_i$ for $1\le i\le N$, the field walk. Its greedy $J$-hysteresis switches first at the first index $g_1$ at which the range $\max_{i'\le i}X_{i'}-\min_{i'\le i}X_{i'}$ reaches $J$. After a switch at which the walk has risen by $J$ above a running minimum (fallen by $J$ below a running maximum), the next phase tracks the running maximum (minimum) from the switch and switches when the walk falls $J$ below it (rises $J$ above it). Let $\Psi_N$ be the number of switches in $\{1,\dots,N\}$.

\begin{lemma}[Counting walls]\label{lem:coh-counting}
% src: L3-L4 (extended to rings, constant adjusted)
If no run of consecutive pulls sums to $\pm J$, then on the ring and on the open chain $\lfloor\Psi_N/2\rfloor-1\le\#\mathrm{walls}(\sigma^\ast)\le2\Psi_N+4$.
\end{lemma}

\begin{proof}
Let $x_1<\dots<x_L$ be the walls at agents $1,\dots,N-1$; on the ring the edge between $N$ and $1$ may carry one more wall. For $1<\alpha\le L$ the agents $x_{\alpha-1}+1,\dots,x_\alpha$ form a domain with walls on both sides, so \cref{lem:coh-structure}(a) gives $\sigma_\alpha(X_{x_\alpha}-X_{x_{\alpha-1}})\ge J$, with $\sigma_\alpha$ its sign. Signs alternate, so the wall levels $X_{x_\alpha}$ alternate between maxima and minima, and at a maximum wall $x_\alpha$ with $1<\alpha<L$, $X_{x_\alpha}\ge X_{x_{\alpha\pm1}}+J$.

Upper bound. Let $x_\alpha$ be a maximum wall with $1<\alpha<L$, and suppose no switch occurs in $(x_{\alpha-1},x_{\alpha+1}]$. If none occurs at or before $x_{\alpha-1}$ either, the range of the walk on $[0,x_\alpha]$ is at least $J$, so $g_1\le x_\alpha$, a contradiction. Otherwise let $g$ be the last switch at or before $x_{\alpha-1}$. If the phase after $g$ tracks a maximum, the running maximum is at least $X_{x_\alpha}$ by $x_\alpha$ and $X_{x_{\alpha+1}}\le X_{x_\alpha}-J$, so a switch occurs after $g$ and by $x_{\alpha+1}$, hence in $(x_{\alpha-1},x_{\alpha+1}]$. If it tracks a minimum, the running minimum is at most $X_{x_{\alpha-1}}$ and $X_{x_\alpha}\ge X_{x_{\alpha-1}}+J$, so a switch occurs in $(x_{\alpha-1},x_\alpha]$. Both contradict the assumption. The windows $(x_{\alpha-1},x_{\alpha+1}]$ of the maximum walls are disjoint, and at least $\lfloor(L-2)/2\rfloor$ maximum walls lie strictly between $x_1$ and $x_L$, so $\lfloor(L-2)/2\rfloor\le\Psi_N$ and $L\le2\Psi_N+3$; the possible wall between $N$ and $1$ adds one.

Lower bound. Between switches $g_k<g_{k+1}<g_{k+2}$ the walk completes reversals of at least $J$ in both directions. If the phase after $g_k$ tracks a maximum, attained at some $i'\in[g_k,g_{k+1})$, then $X_{g_{k+1}}-X_{i'}\le-J$, and the next phase gives $i''\in[g_{k+1},g_{k+2})$ with $X_{g_{k+2}}-X_{i''}\ge J$; the other case is symmetric. The runs of agents $i'+1,\dots,g_{k+1}$ and $i''+1,\dots,g_{k+2}$ lie in $\{g_k+1,\dots,g_{k+2}\}$, and by hypothesis their sums are strictly below $-J$ and strictly above $J$. By \cref{lem:coh-structure}(c) a domain of sign $+1$ contains no run with sum below $-J$ and a domain of sign $-1$ none with sum above $J$, so the agents $g_k+1,\dots,g_{k+2}$ do not lie in one domain, and a wall sits between two consecutive agents of this window. The windows for $k=1,3,5,\dots$ are disjoint, which gives at least $\lfloor(\Psi_N-1)/2\rfloor\ge\lfloor\Psi_N/2\rfloor-1$ walls.
\end{proof}

\begin{lemma}[Drawdown time]\label{lem:coh-drawdown}
% src: L3-L5
Let $(X_u)_{u\ge0}$ be a random walk from $0$ with independent $\Ncal(0,h^2)$ steps, $B_u=\max_{u'\le u}X_{u'}-X_u$ its drawdown and $T_J=\min\{u:B_u\ge J\}$. Then $(J/h)^2\le\E T_J\le13(J/h)^2+3.25$.
\end{lemma}

\begin{proof}
Lower bound. With $h_{u+1}=X_{u+1}-X_u$, $B_{u+1}=\max\{B_u-h_{u+1},0\}$, so $B_{u+1}^2\le(B_u-h_{u+1})^2$ and $\E[B_{u+1}^2\mid X_0,\dots,X_u]\le B_u^2+h^2$; hence $B_u^2-uh^2$ is a supermartingale. Optional stopping at $T_J\wedge n$ gives $\E B^2_{T_J\wedge n}\le h^2\E[T_J\wedge n]$, and letting $n\to\infty$, with Fatou's lemma on the left, monotone convergence on the right and $T_J<\infty$ almost surely by the upper bound, gives $J^2\le\E B_{T_J}^2\le h^2\E T_J$. Upper bound. Split time into blocks of $n_0=\lceil4J^2/h^2\rceil$ steps. A block whose increment is at most $-J$ completes a drawdown of $J$, and this happens independently across blocks with probability $\Pr(\Ncal(0,n_0h^2)\le-J)\ge\Pr(\Ncal(0,1)\le-\frac12)>0.3085$. So $\E T_J\le n_0/0.3085\le(4(J/h)^2+1)/0.3085\le13(J/h)^2+3.25$.
\end{proof}

The lemmas leave a factor of four between the bounds on the number of walls, and the data suggest that it is spurious.

\begin{conjecture}[Exact wall count]\label{conj:coh-walls}
On an open chain the number of domains of the minimizer of the pairwise potential equals the number of switches of the greedy $J$-hysteresis or exceeds it by one, so that the mean domain length tends almost surely to $\E T_J$ of \cref{lem:coh-drawdown}.
\end{conjecture}
% src: L3-C4 (restated through counts: the data compare numbers of domains and switches, not positions)

In exact minimizers on chains of $4\times10^5$ agents the number of domains equaled the number of switches in four runs and exceeded it by one in the fifth (\cref{chk:coh-domains}).

\begin{proof}[Proof of \cref{thm:coh-domains-1d}]
(a) By \cref{prop:coh-code-potential}(b) the game is the pairwise game with coupling $J$ on every edge of the ring, and \cref{prop:coh-code-potential}(a) gives reversibility, the stationary law and its concentration. For distinct configurations $\sigma$ and $\sigma'$, $\Phi_{\mathrm{code}}(\sigma)-\Phi_{\mathrm{code}}(\sigma')$ is a multiple of $J$ plus $-2\sum_{i}h_i\sigma_i$ over the nonempty set where they differ, a variable with a continuous law, so with probability one no two of the finitely many configurations tie and the minimizer is unique.

(b) With probability one no run of pulls sums to $\pm J$, so \cref{lem:coh-counting} applies for every $N$. The switch indices of the greedy hysteresis are stopping times of the walk. After the first, each phase starts at the current value of the walk and ends when the walk has reversed by $J$ from the running extreme, so by the strong Markov property the periods between consecutive switches are independent, and by the symmetry of the steps each has the law of $T_J$ in \cref{lem:coh-drawdown}. Since $\E T_J<\infty$, the renewal strong law of large numbers gives $\Psi_N/N\to1/\E T_J$ almost surely. On the ring the number of domains is the number of walls, or one if there is none, so \cref{lem:coh-counting} gives $N/(2\Psi_N+4)\le\bar\ell_N\le N/(\lfloor\Psi_N/2\rfloor-1)$ once $\Psi_N\ge4$, hence $\frac12\E T_J\le\liminf\bar\ell_N\le\limsup\bar\ell_N\le2\E T_J$ almost surely, and \cref{lem:coh-drawdown} gives the stated bounds, since $2\times3.25<7$.

(c) At $h=0$ the stationary law is proportional to $\exp(\frac{\gamma J}2\sum_i\sigma_i\sigma_{i+1})$, because $\mathbf 1[\sigma_i\ne\sigma_{i+1}]=\frac12(1-\sigma_i\sigma_{i+1})$. The transfer matrix $\mathsf M$ with entries $e^{\gamma J\sigma\sigma'/2}$ has eigenvalues $2\cosh(\gamma J/2)$ and $2\sinh(\gamma J/2)$ with eigenvectors $(1,1)$ and $(1,-1)$, which $\diag(1,-1)$ exchanges. For agents at distance $d$, $\E[\sigma_i\sigma_j]=\tr(\diag(1,-1)\mathsf M^d\diag(1,-1)\mathsf M^{N-d})/\tr\mathsf M^N=(x_\gamma^d+x_\gamma^{N-d})/(1+x_\gamma^N)$ with $x_\gamma=\tanh(\gamma J/2)\in[0,1)$, and $\Pr(\sigma_i\ne\sigma_j)=\frac12(1-\E[\sigma_i\sigma_j])$. The limits follow from $x_\gamma<1$, and $\frac12(1-\tanh(\gamma J/2))=1/(1+e^{\gamma J})$.
\end{proof}

\subsection{The grounding tension}\label{app:coh-tension}

\begin{proof}[Proof of \cref{thm:coh-grounding-tension}]
(a) Since $\hat\Lambda^+>\varpi_{\mathrm s}+2\varpi_{\mathrm y}$, the gains satisfy $k^s+2k^{\mathrm y}<1$, and with $0<F<1$ this gives $c>0$, $a>0$ and $\lambda=F(1-k^s)<1$: \textup{(NN)} and \textup{(C$_r$)} hold. The gains are fixed, so $W$ is deterministic and time invariant, which gives \textup{(Lin)}, and it couples each agent only to itself and its two neighbors, which gives \textup{(Loc)}; the local noises are independent across agents with the constant variance $q$, the world innovation has the constant variance $Q$, and there is no mistranslation term, so the law of the local noise is time invariant and the moment hypotheses of \cref{lem:sca-relations}(a) on the inputs hold, while the stationary law, given by the series of \cref{thm:sca-clustering}(b), involves no initial error. Hence \textup{(M)} follows in stationary form by \cref{thm:sca-clustering}(b). By \cref{lem:fw-common-mode} the stationary error is the sum of a connected part, driven by the local noises, and an independent common mode, which on the ring, where $W\one=\lambda\one$, is
\[ -\sum_{n\ge0}w_{t-1-n}W^n\one=-\sum_{n\ge0}\lambda^nw_{t-1-n}\one, \]
with variance $Q/(1-\lambda^2)=Q/(\varepsilon(2-\varepsilon))$ at every agent. Since $W$ is symmetric, the connected part has covariance $\Sigma^c=q\sum_nW^{2n}=q(I-W^2)^{-1}$, and by Fourier analysis on $\Z$
\[ \Sigma^c_{i,i+d}=\frac q{2\pi}\int_{-\pi}^{\pi}\frac{e^{\mathrm iud}}{1-w(u)^2}\mathrm{d}u,\qquad w(u)=c+2a\cos u. \]
Write $(1-w^2)^{-1}=\frac12[(1-w)^{-1}+(1+w)^{-1}]$ and use, for $A_1>|A_2|>0$, $\frac1{2\pi}\int_{-\pi}^{\pi}e^{\mathrm iud}(A_1-A_2\cos u)^{-1}\mathrm{d}u=(\mathrm{sgn}A_2)^de^{-m_{12}|d|}/\sqrt{A_1^2-A_2^2}$ with $\cosh m_{12}=A_1/|A_2|$. This gives
\[ \Sigma^c_{i,i+d}=\frac q2\Bigl[\frac{e^{-m|d|}}{\sqrt{(1-c)^2-4a^2}}+\frac{(-1)^de^{-m'|d|}}{\sqrt{(1+c)^2-4a^2}}\Bigr],\qquad\cosh m=\frac{1-c}{2a},\quad\cosh m'=\frac{1+c}{2a}, \]
where $1+c>2a$ because $c>0$ and $a\le F/2<\frac12$. At $d=0$, with $(1-c)^2-4a^2=(1-c-2a)(1-c+2a)=\varepsilon(1-c+2a)$, this is the second term of \eqref{eq:coh-ring-mse}, and adding the common mode gives (ii).

(i) Since $c>0$, $m'>m$, so the first term dominates at large $d$ and the standardized connected covariance decays exactly at the rate $m$; hence $\xi=1/m$ by \cref{def:sca-lengths}, with $\cosh m=(1-c)/(2a)=1+\varepsilon/(2a)$, in agreement with \cref{thm:sca-oz}. The gains give $a=F\varpi_{\mathrm y}/\hat\Lambda^+$ and $\varepsilon=1-F+F\varpi_{\mathrm s}/\hat\Lambda^+$, hence the formula for $a/\varepsilon$. As $F\to1$ at fixed precisions $\hat\Lambda^+$ stays bounded, because $\hat P^-\ge Q$, so $(1-F)\hat\Lambda^+\to0$. As $\varpi_{\mathrm y}\to\infty$ at $F<1$, $\hat P^-\ge Q$ keeps $1/\hat P^-$ bounded, so $\hat\Lambda^+/(2\varpi_{\mathrm y})\to1$ and $a/\varepsilon\to F/(2(1-F))$, which makes $1+\varepsilon/(2a)\to1/F$. At every assigned precision $\hat\Lambda^+\ge2\varpi_{\mathrm y}$, so $a/\varepsilon\le F\varpi_{\mathrm y}/(2(1-F)\varpi_{\mathrm y})=F/(2(1-F))$, hence $1+\varepsilon/(2a)\ge1/F$ and $\xi\le1/\arccosh(1/F)$.

(iii) Since $0<\varepsilon<1$, $\varepsilon(2-\varepsilon)\le2\varepsilon=4a(\cosh(1/\xi)-1)=8a\sinh^2(1/(2\xi))$, and $\sinh y\le y\cosh y$ for $y\ge0$ because $\tanh y\le y$. Hence $\varepsilon(2-\varepsilon)\le2a\xi^{-2}\cosh^2(1/(2\xi))$, which gives \eqref{eq:coh-tension-law}. The asymptotic forms follow from $\cosh^2(1/(2\xi))=1+O(\xi^{-2})$ and $a=Fk^{\mathrm y}\le F/2$.

(iv) As $\varpi_{\mathrm y}\to\infty$, $k^s=\varpi_{\mathrm s}/\hat\Lambda^+\to0$ and $k^{\mathrm y}\to\frac12$, so $c\to0$, $a\to F/2$, $\varepsilon\to1-F$ and $q\to F^2R_c/2$. The first term of \eqref{eq:coh-ring-mse} tends to $Q/(1-F^2)=P_0$ and the bracket to $2/\sqrt{1-F^2}$, which gives the limit.

(b) With exact code models and homogeneous sensing gain every row of $W$ sums to $F(1-k)=\lambda$, and by \textup{(NN)} this is also the row-sum norm, so \textup{(C$_r$)} holds. The common mode is $-\sum_n\lambda^nw_{t-1-n}$ at every agent, with variance $Q/(1-\lambda^2)$, independent of the connected part (\cref{lem:fw-common-mode}); this is the bound of \cref{prop:mis-mse-floor}. Under \textup{(Lin)}, \textup{(Loc)}, \textup{(C$_r$)} and \textup{(B)}, \cref{thm:sca-clustering} gives $\xi\le1/\ln(1/\lambda)$, so $\lambda^2\ge e^{-2/\xi}$ and $\E e_i^2\ge Q/(1-\lambda^2)\ge Q/(1-e^{-2/\xi})$, and $1-e^{-y}\le y$ gives the last inequality.
\end{proof}

\begin{proof}[Proof of \cref{prop:coh-placement}]
(a) The receiver's predictive density of a neighbor's message is $\Ncal(m_i,S)$ with $S=\hat P^-+1/\varpi_{\mathrm y}$. At the exact posterior of its model the variational free energy of one message equals its negative log evidence, $\frac12\ln(2\pi S)+(y_{ij}-m_i)^2/(2S)$, so the time average is $\frac12\ln S$ plus the time average of $(y_{ij}-m_i)^2$ divided by $2S$, up to a constant; in stationarity that time average is $\E(y_{ij}-m_i)^2$. As a function of $S$ this decreases up to $S=\E(y_{ij}-m_i)^2$ and increases beyond, and the constraint $\varpi_{\mathrm y}\le1/R_c$ is $S\ge\hat P^-+R_c$, which gives (a). (b) $y_{ij}-m_i=e_j-e_i+\nu_{ij}$. On the ring the common mode is the same at every agent, so $e_j-e_i$ is the difference of the connected parts, with second moment $2(\Sigma^c_{ii}-\Sigma^c_{ij})$ by translation invariance, and the channel noise adds $R_c$. Under the stated condition the target $\E(y_{ij}-m_i)^2-\hat P^-$ of (a) is at most $R_c$. (c) $s_i-m_i=v_i-e_i$ with $v_i$ independent of the prior error $e_i$, so $\E(s_i-m_i)^2=R^s+\E e_i^2$. The receiver's model gives this innovation the variance $\hat P^-+\hat R^s$, and moment matching sets the assigned sensing noise $\hat R^s=R^s+\E e_i^2-\hat P^-$, which exceeds $R^s$ exactly when $\E e_i^2>\hat P^-$.
\end{proof}
